% !PRINTED 578
\weberpage{578}

%<<CONT
can set. Now since in the sequence of numbers $\nu^2$ the number 0
occurs, and moreover every quadratic residue occurs twice [for
$\nu^2$ and $(-\nu)^2$], we have

\[ r^{\frac{m^2}{4}} f(r^{-2}, m) = 1 + 2A = A - B, \]

and from formula (12) it follows that

\[ A - B = r^{\frac{m^2}{4}} (1 - r^{-2})(1 - r^{-6}) \ldots (1 - r^{-2n + 4}). \]

Now

\[ r^{\frac{m^2}{4}} = r\, r^3 r^5 \ldots r^{n - 2}, \]

and so we finally obtain the formula

\weq{13}{A - B = (r - r^{-1})(r^3 - r^{-3}) \ldots (r^{n - 2} - r^{-n + 2}).}

If $\nu$ denotes an odd and $\mu$ an even number between
0 and ${}^1/_2 n$, this product consists of the factors

\[ r^{\nu} - r^{-\nu}, \quad r^{n - \mu} - r^{-n + \mu} = -(r^{\mu} - r^{-\mu}). \]

The number of factors of the second category is as great
as the number of even numbers between 0 and ${}^1/_2 n$, or
equal to the number of positive integers that are smaller than
${}^1/_4 n$. If therefore, understanding by $l$ an integer, we set
$n = 4l + 1$ or $4l + 3$, then $l$ is this number, and
we can set

\[ A - B = (-1)^{l} \overset{\nu}{\Pi} (r^{\nu} - r^{-\nu}) \overset{\mu}{\Pi} (r^{\mu} - r^{-\mu}), \]

for which, if $h$ runs through the sequence of numbers from 1 to ${}^1/_2 (n - 1)$,
we may also set

\[ A - B = (-1)^{l} \overset{h}{\underset{1,\, \frac{n - 1}{2}}{\Pi}} (r^{h} - r^{-h}). \]

If we now set

\[ r = e^{\frac{2 \pi i k}{n}}, \quad r^{h} - r^{-h} = 2i \sin \frac{2 \pi k h}{n}, \]

then

\[ A - B = (-1)^{l} (2i)^{\frac{n - 1}{2}} \overset{h}{\Pi} \sin \frac{2 \pi h k}{n}. \]

The product that still occurs here we have already
determined earlier [\S.\,138, (4)]:

\[ 2^{\frac{n - 1}{2}} \overset{h}{\Pi} \sin \frac{2 \pi k h}{n} = \left( \frac{k}{n} \right) \sqrt{n}, \]

wherein the $\sqrt{n}$ is to be taken positive.
